\section{ROS2 C++ Style Guide}
\label{sec:ros2_cpp_style}

\subsection{Overview}

The STAR project uses different coding conventions for ROS2 C++ code compared to RX72N firmware. This section documents the ROS2 C++ style guide, which supplements the C firmware style guide (Section~\ref{sec:c-style}).

\subsubsection{When to Use This Guide
}

\begin{itemize}
  \item \textbf{ROS2 packages} (\texttt{
  star - ros2 / src /*/}): Use ROS2 C++ style
\item \textbf{RX72N firmware} (\texttt{star-rx72n-firmware/}): Use C firmware
style \item \textbf{Gateway} (\texttt{star-gateway/}): Follow Go conventions
\end{itemize}

\subsubsection{Key Differences}

\begin{table}[H]
\centering
\begin{tabular}{|l|l|l|}
\hline
\textbf{Feature} & \textbf{C Firmware (RX72N)} & \textbf{ROS2 C++} \\
\hline
Classes & N/A & \texttt{CamelCase} \\
Methods & \texttt{snake\_case} & \texttt{snake\_case} (same) \\
Variables & \texttt{snake\_case} & \texttt{snake\_case} (same) \\
Member vars & \texttt{s\_} prefix & trailing \texttt{\_} \\
Headers & \texttt{.h} & \texttt{.hpp} \\
Line limit & 100 characters & 120 characters \\
Namespaces & Not used & \texttt{star::package\_name::} \\
Error handling & Return codes & Exceptions \\
Logging & \texttt{uart\_puts()} & \texttt{RCLCPP\_INFO/WARN/ERROR} \\
Include guards & \texttt{STAR\_RX72N\_FILE\_H} & \texttt{PACKAGE\_FILE\_HPP\_}
\\
Doxygen & \texttt{/**} and \texttt{/**<} & \texttt{/**} and \texttt{/**<} (same)
\\ \hline \end{tabular} \caption{C vs C++ Style Comparison}
\label{tab:c_vs_cpp_style}
\end{table}

\subsection{Naming Conventions}

\subsubsection{Classes and Types}

Use \texttt{CamelCase} for classes following ROS2 conventions:

\begin{lstlisting}[language=C++, caption={Class naming examples},
label={lst:class_naming}]
// CamelCase for classes (ROS2 convention)
class StarGatewayBridgeNode : public rclcpp::Node {
// ...
};

// CamelCase for structs used as types
struct TelemetryData {
double battery_voltage_;
int32_t current_ma_;
};

// Type aliases use CamelCase
using TelemetryPtr = std::shared_ptr<TelemetryData>;
\end{lstlisting}

\subsubsection{Methods and Functions}

Use \texttt{snake\_case} for method names (same as C firmware):

\begin{lstlisting}[language=C++, caption={Method naming examples},
label={lst:method_naming}]
// snake_case for methods (same as C firmware)
void publish_telemetry(const TelemetryData & data);
bool is_connected() const;
void on_battery_state_received(const sensor_msgs::msg::BatteryState::SharedPtr
msg);

// Use verb-based names that clarify actions
void check_for_errors();      // Good: Clear intent
void error_check();           // Bad: Noun-first is confusing
\end{lstlisting}

\subsubsection{Variables}

Use \texttt{under\_scored} for variables with trailing underscores for member
variables:

\begin{lstlisting}[language=C++, caption={Variable naming examples},
label={lst:variable_naming}]
// under_scored for variables
int loop_counter = 0;
std::string node_name = "gateway_bridge";

// Member variables with trailing underscore
class MyNode : public rclcpp::Node {
private:
rclcpp::Publisher<std_msgs::msg::String>::SharedPtr telemetry_pub_;
std::shared_ptr<grpc::Channel> grpc_channel_;
bool is_connected_;
};

// Constants: ALL_CAPITALS
const int MAX_RETRIES = 3;
const double DEFAULT_TIMEOUT_S = 5.0;
\end{lstlisting}

\subsubsection{Namespaces}

Package-based namespaces use \texttt{under\_scored} naming:

\begin{lstlisting}[language=C++, caption={Namespace examples},
label={lst:namespace}]
namespace star { namespace spi_bridge {

class SpiDriverNode : public rclcpp::Node {
// ...
};

}  // namespace spi_bridge
}  // namespace star
\end{lstlisting}

\subsection{File Naming and Organization}

\subsubsection{Header Files}

C++ headers use the \texttt{.hpp} extension (not \texttt{.h}):

\begin{itemize}
\item \texttt{star\_gateway\_bridge\_node.hpp}
\item \texttt{message\_converter.hpp}
\item \texttt{grpc\_client.hpp}
\end{itemize}

Include guards follow the pattern: \texttt{PACKAGE\_FILE\_NAME\_HPP\_}

\begin{lstlisting}[language=C++, caption={Include guard example},
label={lst:include_guard}]
#ifndef STAR_GATEWAY_BRIDGE_STAR_GATEWAY_BRIDGE_NODE_HPP_
#define STAR_GATEWAY_BRIDGE_STAR_GATEWAY_BRIDGE_NODE_HPP_

// ... code ...

#endif  // STAR_GATEWAY_BRIDGE_STAR_GATEWAY_BRIDGE_NODE_HPP_
\end{lstlisting}

\subsubsection{Source Files}

Source files use the \texttt{.cpp} extension with \texttt{under\_scored} naming:

\begin{itemize}
\item \texttt{star\_gateway\_bridge\_node.cpp}
\item \texttt{message\_converter.cpp}
\item \texttt{grpc\_client.cpp}
\end{itemize}

\subsection{Header Organization}

Headers must be organized in a standard order (enforced by
\texttt{ament\_cpplint}, ROS2's official include-order checker):

\begin{enumerate}
\item License and copyright
\item Include guard
\item ROS2 core includes (\texttt{<rclcpp/...>})
\item ROS2 message includes (\texttt{<geometry\_msgs/...>}, etc.)
\item Project includes (\texttt{"star/...})
\item System C++ includes (\texttt{<memory>}, \texttt{<string>}, etc.)
\item System C includes (\texttt{<cstdint>}, etc.)
\item Namespace declaration
\item Class/function declarations
\end{enumerate}

See CLAUDE.md for complete example.

\subsection{Code Formatting}

\subsubsection{Line Length}

\begin{itemize}
\item \textbf{ROS2 C++}: 100 characters (enforced by
\texttt{ament\_uncrustify}, ROS2's official formatter)
\item \textbf{C firmware}: 100 characters
(\texttt{star-rx72n-firmware/.clang-format}) \end{itemize}

\textbf{Note:} The ROS2 packages do NOT carry a \texttt{star-ros2/.clang-format}
file. ROS2 uses \texttt{ament\_uncrustify} (uncrustify with a bundled
\texttt{ament\_code\_style.cfg}) and \texttt{ament\_cpplint}, NOT clang-format.
Running clang-format on \texttt{star-ros2/} sources will introduce drift that
\texttt{ament\_uncrustify --check} rejects in CI. To format ROS2 code locally,
source ROS2 and run \texttt{./scripts/ros2/format-ros2.sh}.

\subsubsection{Indentation}

All code uses 2-space indentation (same as C firmware):

\begin{lstlisting}[language=C++, caption={Indentation example},
label={lst:indentation}]
class MyNode : public rclcpp::Node { public: MyNode() :
Node("my_node") { timer_ = this->create_wall_timer(
std::chrono::milliseconds(100),
std::bind(&MyNode::timerCallback, this));
}

private:
void timerCallback() {
// ...
}

rclcpp::TimerBase::SharedPtr timer_;
};
\end{lstlisting}

\subsubsection{Braces}

\begin{itemize}
\item \textbf{Functions}: Braces on new line (same as C firmware)
\item \textbf{Control statements}: Cuddled braces (\texttt{if (x) \{})
\item \textbf{Namespaces}: No indentation inside namespace
\end{itemize}

\begin{lstlisting}[language=C++, caption={Brace style examples},
label={lst:braces}]
// Functions: Braces on new line
void myFunction()
{
// ...
}

// Control statements: Cuddled braces
if (condition) {
// ...
} else {
// ...
}

// Namespaces: No indentation inside
namespace star {
namespace spi_bridge {

// Content at zero indent
class MyClass {
};

}  // namespace spi_bridge
}  // namespace star
\end{lstlisting}

\subsection{ROS2-Specific Patterns}

\subsubsection{Node Inheritance}

Inherit from \texttt{rclcpp::Node} for basic nodes or
\texttt{rclcpp\_lifecycle::LifecycleNode} for safety-critical nodes:

\begin{lstlisting}[language=C++, caption={Node inheritance},
label={lst:node_inheritance}]
// Basic node
class StarGatewayBridgeNode : public rclcpp::Node {
public:
StarGatewayBridgeNode();

private:
void telemetryCallback();

rclcpp::Publisher<std_msgs::msg::String>::SharedPtr telemetry_pub_;
rclcpp::TimerBase::SharedPtr telemetry_timer_;
};

// Safety-critical lifecycle node
class StarSpiDriverNode : public rclcpp_lifecycle::LifecycleNode {
public:
using CallbackReturn =
rclcpp_lifecycle::node_interfaces::LifecycleNodeInterface::CallbackReturn;

StarSpiDriverNode();

// Lifecycle transitions
CallbackReturn on_configure(const rclcpp_lifecycle::State &) override;
CallbackReturn on_activate(const rclcpp_lifecycle::State &) override;
CallbackReturn on_deactivate(const rclcpp_lifecycle::State &) override;
CallbackReturn on_cleanup(const rclcpp_lifecycle::State &) override;
};
\end{lstlisting}

\subsubsection{Publishers and Subscribers}

\begin{lstlisting}[language=C++, caption={Publishers and subscribers},
label={lst:pub_sub}]
class MyNode : public rclcpp::Node { public: MyNode() :
Node("my_node") {
// Publisher: use SharedPtr
telemetry_pub_ = this->create_publisher<std_msgs::msg::String>(
"/telemetry",
10  // QoS depth
);

// Subscriber: use std::bind
cmd_sub_ = this->create_subscription<geometry_msgs::msg::Twist>(
"/cmd_vel",
10,
std::bind(&MyNode::cmdVelCallback, this, std::placeholders::_1)
);
}

private:
void cmdVelCallback(const geometry_msgs::msg::Twist::SharedPtr msg) {
// Handle message
}

rclcpp::Publisher<std_msgs::msg::String>::SharedPtr telemetry_pub_;
rclcpp::Subscription<geometry_msgs::msg::Twist>::SharedPtr cmd_sub_;
};
\end{lstlisting}

\subsubsection{Timers}

Use \texttt{create\_wall\_timer} for periodic operations:

\begin{lstlisting}[language=C++, caption={Timer usage}, label={lst:timers}]
timer_ = this->create_wall_timer(
std::chrono::milliseconds(100),  // 10 Hz
std::bind(&MyNode::timerCallback, this)
);
\end{lstlisting}

\subsection{Error Handling}

ROS2 uses exceptions (different from C firmware return codes):

\begin{lstlisting}[language=C++, caption={Exception handling},
label={lst:exceptions}]
// Throw exceptions for errors
void publishTelemetry()
{
if (!grpc_channel_) {
throw std::runtime_error("gRPC channel not initialized");
}
// ...
}

// Catch exceptions in callbacks (avoid crashing node)
void timerCallback()
{
try {
publishTelemetry();
} catch (const std::exception & e) {
RCLCPP_ERROR(this->get_logger(), "Failed to publish: %s", e.what());
}
}
\end{lstlisting}

\subsection{Logging}

Use \texttt{RCLCPP\_*} macros for logging, not \texttt{printf} or \texttt{cout}:

\begin{lstlisting}[language=C++, caption={ROS2 logging}, label={lst:logging}]
// ROS2 logging macros (preferred)
RCLCPP_INFO(this->get_logger(), "Node started");
RCLCPP_WARN(this->get_logger(), "Connection lost, retrying...");
RCLCPP_ERROR(this->get_logger(), "Failed to initialize: %s", error_msg.c_str());
RCLCPP_DEBUG(this->get_logger(), "Processing message %d", count);

// Throttled logging (max once per 5 seconds)
RCLCPP_WARN_THROTTLE(this->get_logger(), *this->get_clock(), 5000,
"Command stale (%ldms > %dms)", cmd_age_ms, timeout_ms_);
\end{lstlisting}

\textbf{Important}: Never use \texttt{printf()} or \texttt{std::cout} in ROS2
nodes.

\subsection{Documentation}

All public APIs must have Doxygen comments using \texttt{/**} and \texttt{/**<}
(same as C firmware):

\begin{lstlisting}[language=C++, caption={Doxygen documentation},
label={lst:doxygen}]
/**
* @brief Brief description of class
*
* Detailed description can span multiple lines. Explain purpose,
* usage, and any important details.
*/
             class StarGatewayBridgeNode : public rclcpp::Node {
  public:
    /**
     * @brief Constructor for gateway bridge node
     *
     * @param options Node options for configuration
     */
    explicit StarGatewayBridgeNode(
        const rclcpp::NodeOptions &options = rclcpp::NodeOptions());

  private:
    /**
     * @brief Callback for telemetry publishing timer
     *
     * Forwards latest telemetry to Gateway via gRPC. Called at 10 Hz.
     */
    void telemetry_timer_callback();

    rclcpp::TimerBase::SharedPtr telemetry_timer_; /**< Timer for telemetry */
 };
  \end{lstlisting}

  \subsection{Constants}

  Prefer enums and \texttt{const}(same philosophy as C firmware)
      :

\begin{lstlisting}[language = C++, caption = {Constants},
                    label = {lst : constants}]
        // Enum for related constants
        enum class MotorState {
          kIdle = 0,
          kRunning = 1,
          kError = 2
       };

  // const for single values
  const int DEFAULT_QOS_DEPTH = 10;
  const double MAX_LINEAR_VELOCITY_MPS = 1.0;

  // static const for class-specific constants
  class MyNode : public rclcpp::Node {
  private:
    static constexpr int kMaxRetries = 3;
    static constexpr double kTimeoutS = 5.0;
 };
  \end{lstlisting}

  \subsection{Formatting Enforcement}

  All ROS2 C++ code uses \texttt{clang - format} with configuration in \texttt{
      star - ros2 /.clang - format} :

\begin{itemize}
  \item 2-space indentation (same as C firmware)
  \item 120-character line limit (ROS2 standard)
  \item Cuddled braces for control statements
  \item Function braces on new line
  \item No indentation inside namespaces
  \item ROS2-specific include order (core, messages, system, project)
\end{itemize}

  \subsubsection{Manual Formatting}

          Format code before committing :

\begin{verbatim} cd star -
      ros2 find src - name '*.cpp' - o - name '*.hpp' |
      xargs clang - format - i
\end{verbatim}

  \subsubsection{CI Enforcement}

          The \texttt{.github / workflows / ros2.yml} workflow runs \texttt{
              clang - format} checks on all PRs :

\begin{verbatim} find src -
      name '*.cpp' - o - name '*.hpp' |
      xargs clang - format-- dry -
          run-- Werror
\end{verbatim}

          If formatting check fails,
      the build will fail with instructions to fix.

\subsection{References
 }

  \begin{itemize}
  \item \textbf{ROS2 C++ Style Guide} : \url {
  https: // docs.ros.org/en/rolling/The-ROS2-Project/Contributing/Code-Style-Language-Versions.html}
    \item \textbf{ROS C++ Best Practices} : \url {
    https: // wiki.ros.org/CppStyleGuide}
      \item \textbf{Google C++ Style Guide} : \url {
      https: // google.github.io/styleguide/cppguide.html}
        \item \textbf{STAR C Firmware Style} : CLAUDE.md Section
                                               "Code Style"
  \item \textbf{STAR Project Documentation} : \texttt {
          docs / star\_documentation.pdf
       }
        \end{itemize}

        \subsection{Integration with Existing Tools}

        The ROS2 C++ style guide integrates seamlessly with existing STAR
            project tools :

\begin{table}[H]
\centering
\begin{tabular} {
          | l | l | l |
       }
        \hline
\textbf{Tool} & \textbf{Purpose} & \textbf{Configuration}
        \\
\hline
\texttt{clang - format} & Code formatting & \texttt {
          star - ros2 /.clang - format
       }
        \\
\texttt{ament\_cppcheck} & Static analysis &ROS2 workflow \\
\texttt{ament\_cpplint} & Style checking &ROS2 workflow \\
\texttt{colcon build} & Build system &CMakeLists.txt \\
\texttt{colcon test} &
            Unit testing &gtest integration \\
GitHub Actions &CI /
                CD enforcement & \texttt {
          .github / workflows / ros2.yml
       }
        \\
\hline
\end{tabular}
        \caption{ROS2 Development Tools}
        \label{tab:ros2_tools}
        \end{table}

\subsection{Summary}

The STAR project maintains separate but complementary style guides for C firmware and C++ ROS2 code:

\begin{itemize}
\item \textbf{C firmware} follows embedded best practices with strict NASA Power
    of 10 compliance
  \item \textbf{ROS2 C++} follows ROS community
        conventions while maintaining STAR project philosophy
  \item Both styles share common principles
    : enums over macros,
      error checking,
      documentation
  \item Automated tooling enforces consistency across all code
\end{itemize}

      For complete examples and additional details,
      refer to CLAUDE.md.
